\chapter{Sistema de Entidades}
\label{ch:entities}

\section{Modelo orientado a objetos}

Toda forma de vida (e morte) no mundo do Void Descent é uma \texttt{entity\_t}, uma
classe-base ECMAScript 2015 com herança simples. O \textbf{Class Diagram} comunica
simultaneamente herança, atributos, métodos e relacionamentos.

\diagfiglandscape{diagrams/entities-class.png}{Diagrama de Classes, com a hierarquia de \texttt{entity\_t}}{entities-class}

\section{A classe-base \texttt{entity\_t}}

Todas as entidades compartilham:

\begin{itemize}
  \item \textbf{Cinemática}: posição (\texttt{x, y, z}), velocidade
    (\texttt{vx, vy, vz}), aceleração (\texttt{ax, ay, az}) e atrito (\texttt{f}).
  \item \textbf{Visualização}: \texttt{s} (índice do tile no atlas) e \texttt{\_dead}
    (\emph{flag} de remoção).
  \item \textbf{Saúde}: \texttt{h} (\emph{hit points}), inicializado a 5 e
    sobrescrevível no \texttt{\_init}.
\end{itemize}

A interface comportamental é definida por seis métodos sobrescrevíveis:

\begin{description}
  \item[\texttt{\_init(param)}] \emph{Hook} de inicialização chamado pelo construtor.
    Permite à subclasse personalizar atributos sem reimplementar a cadeia de
    inicialização.
  \item[\texttt{\_update()}] Avança a física por \texttt{time\_elapsed} segundos. A
    implementação base aplica forças, integra a posição e testa colisão com paredes via
    \texttt{\_collides}.
  \item[\texttt{\_collides(x, z)}] Checa se a posição alvo está dentro de um tile
    sólido. Sobrescrito em \texttt{entity\_player\_t} para implementar a barreira do
    \emph{lock-and-clear}.
  \item[\texttt{\_check(other)}] Chamado para cada par de entidades cuja AABB se
    intersecta. Implementa \emph{pickup}, dano por contato, separação entre inimigos,
    entre outros.
  \item[\texttt{\_render()}] Emite primitivas para o \emph{renderer}
    (\texttt{push\_sprite}, \texttt{push\_light}).
  \item[\texttt{\_kill()}] Marca a entidade como \texttt{\_dead}, agendando sua remoção
    no fim do \emph{frame}.
\end{description}

\section{Subclasses de combate}

\subsection{\texttt{entity\_player\_t}}

O jogador é único em vários aspectos:

\begin{itemize}
  \item Sobrescreve \texttt{\_collides} para impedir saída de salas trancadas.
  \item Mantém estado de \emph{combo} (\texttt{combo}, \texttt{combo\_timer}) e
    pontuação.
  \item Gerencia o \emph{cooldown} do \emph{dash} (\texttt{\_dash\_cooldown},
    \texttt{\_dashing}) e da arma equipada (\texttt{\_last\_shot}).
  \item Em \texttt{\_update}, calcula o ângulo de mira a partir da posição do mouse e
    do \emph{offset} da câmera, e despacha o ataque pelo \texttt{type} da arma
    (\texttt{melee} ou \texttt{ranged}).
\end{itemize}

\subsection{Inimigos}

\begin{description}
  \item[\texttt{entity\_seeker\_t}] (HP 3). Persegue o jogador em linha reta com
    \emph{retarget} periódico. Comportamento mais simples, introduzido no Andar 1.
  \item[\texttt{entity\_shooter\_t}] (HP 5). Mantém distância de 24 unidades do
    jogador, e dispara em cone duas \texttt{entity\_enemy\_plasma\_t} ao detectá-lo.
    Aparece a partir do Andar 2.
  \item[\texttt{entity\_bomber\_t}] (HP 2). Persegue agressivamente; ao chegar a 20
    unidades de distância, ativa um \emph{fuse} de 0,8\,s e explode, causando dano em
    raio de 28 unidades. Introduzido no Andar 3.
  \item[\texttt{entity\_boss\_t}] (HP 70). Adversário do Andar 3, com duas fases:
    \begin{enumerate}
      \item Fase 1 (HP > 35): atira anéis de 7 plasmas em padrão rotativo.
      \item Fase 2 (HP $\leq$ 35): abandona o padrão rotativo, avança em \emph{melee}
        e invoca grupos de 3 \texttt{entity\_seeker\_t}.
    \end{enumerate}
\end{description}

\section{Subclasses de projétil e efeito}

\begin{description}
  \item[\texttt{entity\_plasma\_t}] Projétil do jogador (\emph{ranged}). Move-se em
    linha reta, hospeda \texttt{\_dmg} configurável pela arma, e é destruído ao colidir
    com parede ou inimigo.
  \item[\texttt{entity\_enemy\_plasma\_t}] Versão adversária, que só causa dano ao
    jogador.
  \item[\texttt{entity\_explosion\_t}] Efeito visual transitório de 1 segundo, com luz
    pulsante para \emph{feedback} de impacto.
  \item[\texttt{entity\_slash\_t}] \emph{Burst} de partículas com \emph{flash} de luz
    para o golpe \emph{melee}. \emph{Lifetime} curto, de 180\,ms.
  \item[\texttt{entity\_particle\_t}] Partícula com gravidade e \emph{bounce}, gerada
    por impactos.
\end{description}

\section{Subclasses de pickup e cenário}

\begin{description}
  \item[\texttt{entity\_health\_t}] \emph{Pickup} de poção (curativo de 1 HP). Adiciona
    ao \emph{slot} de consumível do jogador (capacidade 1). Spawn aleatório, com 50\%
    de chance por sala regular, ou garantido na sala do baú.
  \item[\texttt{entity\_weapon\_t}] \emph{Pickup} de arma. Sorteia um índice diferente
    da arma atualmente equipada e substitui ao toque.
  \item[\texttt{entity\_staircase\_t}] Escada de descida. Inicialmente inativa
    (\texttt{\_active = false}), é ativada por \texttt{count\_enemies} quando todos os
    inimigos da sala estão mortos. Ao ser tocada, dispara \texttt{show\_shop} (em
    andares não-finais) ou termina o jogo (no último andar).
\end{description}

\section{Lifecycle e remoção}

A remoção de entidades segue um \textit{deferred deletion} para preservar a
estabilidade da iteração:

\begin{enumerate}
  \item \texttt{\_kill()} marca \texttt{\_dead = true} e empurra a referência em
    \texttt{entities\_to\_kill}.
  \item Iterações subsequentes do \texttt{\_update}, \texttt{\_check} e
    \texttt{\_render} no mesmo \emph{frame} testam \texttt{\_dead} e pulam a entidade.
  \item No fim do \emph{frame}, \texttt{entities = entities.filter(\dots)} reconstrói
    o \emph{array} sem as entidades marcadas.
\end{enumerate}
